% =================================================
% Ученический рабочий лист. Урок 3
% =================================================

\worksheettitle
  {Составление чисел. Задачи на цифры}
  {Урок 3 из 3. Обозначение натуральных чисел}

\studentnameblock

\begin{checkpointbox}[Входная диагностика: 4 минуты]
  \begin{enumerate}[label=\textbf{\arabic*.},itemsep=0.38em]
    \item Какая цифра стоит в разряде десятков числа \(7\,405\)?
      Ответ: \answerblank[14mm].
    \item Запишите цифрами: пять миллионов восемь тысяч два.
      \answerblank[40mm].
    \item Составьте наибольшее трёхзначное число из цифр \(1,4,7\),
      используя каждую один раз: \answerblank[20mm].
    \item Является ли запись \(053\) трёхзначным числом?
      \choicebox\ да \quad \choicebox\ нет.
  \end{enumerate}
\end{checkpointbox}

\section{Порядок цифр изменяет число}

\orderchangesfigure

Заполните:
\[
  27=\answerblank[12mm]\text{ десятка}+\answerblank[12mm]\text{ единиц},
\]
\[
  72=\answerblank[12mm]\text{ десятков}+\answerblank[12mm]\text{ единицы}.
\]

\begin{reflectionbox}[Сделайте вывод]
  Числа \(27\) и \(72\) составлены из одинаковых цифр, но не равны, потому
  что \answerblank[75mm].
\end{reflectionbox}

\section{Цифры могут повторяться}

Нужно записать все двузначные числа, составленные только из цифр \(3\) и
\(7\). Повторение цифр разрешено.

\repeatstreefigure

\begin{fillbox}[Заполните ветви перебора]
  \begin{enumerate}
    \item Если цифра десятков равна \(3\), получаем числа
      \answerblank[15mm] и \answerblank[15mm].
    \item Если цифра десятков равна \(7\), получаем числа
      \answerblank[15mm] и \answerblank[15mm].
    \item Всего получено \answerblank[12mm] числа.
  \end{enumerate}
\end{fillbox}

\begin{thinkbox}[Почему список полный?]
  Объясните, почему другого двузначного числа из этих цифр получить нельзя.
  \writinglines{2}
\end{thinkbox}

\begin{practicebox}[Попробуйте по тому же плану]
  Запишите все двузначные числа из цифр \(2\) и \(5\), если цифры могут
  повторяться:
  \[
    \answerblank[72mm].
  \]
\end{practicebox}

\section{Цифры не повторяются}

Нужно составить все трёхзначные числа из цифр \(0,5,8\), используя каждую
цифру один раз.

\norepeattreefigure

\begin{fillbox}[Закончите перебор]
  \begin{enumerate}
    \item Первая цифра не может быть равна \answerblank[12mm].
    \item Если первая цифра \(5\), получаем числа
      \answerblank[18mm] и \answerblank[18mm].
    \item Если первая цифра \(8\), получаем числа
      \answerblank[18mm] и \answerblank[18mm].
  \end{enumerate}
\end{fillbox}

\begin{warningbox}[Начальный ноль]
  Дополните:
  \[
    053=\answerblank[18mm].
  \]
  Поэтому запись \(053\) является \answerblank[28mm] числом, а не
  трёхзначным.
\end{warningbox}

\begin{practicebox}[Без повторений]
  Запишите все трёхзначные числа из цифр \(0,4,9\), если каждая цифра
  используется один раз:
  \writinglines{1}
\end{practicebox}

\section{Наибольшее и наименьшее число}

\largestsmallestfigure

\begin{fillbox}[Сформулируйте правила]
  \begin{enumerate}
    \item Для наибольшего числа цифры располагают по
      \answerblank[28mm].
    \item Для наименьшего числа без нуля цифры располагают по
      \answerblank[28mm].
    \item Если есть ноль, на первое место ставят
      \answerblank[48mm].
  \end{enumerate}
\end{fillbox}

Заполните таблицу.

\begin{center}
\renewcommand{\arraystretch}{1.55}
\begin{tabularx}{\linewidth}{@{}>{\centering\arraybackslash}p{0.28\linewidth}>{\centering\arraybackslash}X>{\centering\arraybackslash}X@{}}
  \toprule
  \rowcolor{TableHeader}
  Цифры & Наибольшее число & Наименьшее число \\
  \midrule
  \(1,3,6,9\) & \answerblank[25mm] & \answerblank[25mm] \\
  \(0,2,4,7\) & \answerblank[25mm] & \answerblank[25mm] \\
  \(0,1,5,8,9\) & \answerblank[30mm] & \answerblank[30mm] \\
  \bottomrule
\end{tabularx}
\end{center}

\section{Условие связывает цифры}

Найдите все двузначные числа, в которых цифра десятков на \(3\) больше
цифры единиц.

\systematicroutefigure

Заполните таблицу.

\begin{center}
\renewcommand{\arraystretch}{1.45}
\begin{tabular}{c|ccccccc}
  \toprule
  цифра единиц & 0 & 1 & 2 & 3 & 4 & 5 & 6 \\
  \midrule
  цифра десятков & \answerblank[7mm] & \answerblank[7mm] & \answerblank[7mm] &
  \answerblank[7mm] & \answerblank[7mm] & \answerblank[7mm] & \answerblank[7mm] \\
  \bottomrule
\end{tabular}
\end{center}

Ответ:
\[
  \answerblank[105mm].
\]

\begin{thinkbox}[Почему перебор заканчивается?]
  Почему нельзя взять цифру единиц \(7\)?
  \writinglines{2}
\end{thinkbox}

\begin{practicebox}[Новое условие]
  Запишите все двузначные числа, в которых цифра единиц в \(2\) раза больше
  цифры десятков.

  \begin{center}
  \renewcommand{\arraystretch}{1.35}
  \begin{tabular}{c|cccc}
    \toprule
    цифра десятков & 1 & 2 & 3 & 4 \\
    \midrule
    цифра единиц & \answerblank[8mm] & \answerblank[8mm] &
      \answerblank[8mm] & \answerblank[8mm] \\
    \bottomrule
  \end{tabular}
  \end{center}

  Ответ: \answerblank[70mm].
\end{practicebox}

\section{Считаем подходящие числа}

Сколько существует двузначных чисел, сумма цифр которых равна \(7\)?

\textbf{Заполните пары цифр.}

\begin{center}
\renewcommand{\arraystretch}{1.22}
\begin{tabular}{c|ccccccc}
  \toprule
  десятки & 1 & 2 & 3 & 4 & 5 & 6 & 7 \\
  \midrule
  единицы & \answerblank[7mm] & \answerblank[7mm] & \answerblank[7mm] &
  \answerblank[7mm] & \answerblank[7mm] & \answerblank[7mm] & \answerblank[7mm] \\
  \bottomrule
\end{tabular}
\end{center}

\noindent Числа: \answerblank[92mm].\qquad
Количество: \answerblank[15mm].

\section{Самостоятельная проверка}

\begin{practicebox}[Выполните без подсказок]
  \begin{enumerate}[label=\textbf{\arabic*.},itemsep=0.45em]
    \item Запишите все двузначные числа из цифр \(4\) и \(9\), если цифры
      могут повторяться: \answerblank[72mm].
    \item Запишите все трёхзначные числа из цифр \(0,2,7\) без повторений:
      \answerblank[82mm].
    \item Из цифр \(0,3,6,8\) составьте наибольшее и наименьшее
      четырёхзначные числа: \answerblank[22mm] и \answerblank[22mm].
    \item Запишите все двузначные числа, в которых цифра десятков на \(2\)
      больше цифры единиц:
      \writinglines{1}
    \item Какое число получится из \(5\) десятков тысяч и \(78\) сотен?
      \answerblank[34mm].
  \end{enumerate}
\end{practicebox}

\begin{checkpointbox}[Выходной билет]
  \begin{enumerate}
    \item Почему \(041\) не является трёхзначным числом?
      \writinglines{1}
    \item Наименьшее четырёхзначное число из цифр \(0,2,5,9\):
      \answerblank[24mm].
    \item Число, у которого цифра десятков \(6\), а цифра единиц на \(4\)
      меньше: \answerblank[18mm].
  \end{enumerate}
\end{checkpointbox}

\begin{reflectionbox}[Оцените себя]
  \choicebox\ могу объяснить полный перебор;
  \quad
  \choicebox\ умею по образцу;
  \quad
  \choicebox\ нужна тренировка.
\end{reflectionbox}
